% Standalone problem document, generated from the adjacent README.md.
% Compile with XeLaTeX. Mathematical content is maintained in Markdown.
\documentclass[11pt,a4paper]{article}
\usepackage[fontsize=10.5pt]{scrextend}
\usepackage[margin=22mm,headheight=15pt,headsep=9mm,footskip=11mm]{geometry}
\usepackage{fontspec}
\setmainfont{texgyrepagella}[Extension=.otf,UprightFont=*-regular,BoldFont=*-bold,ItalicFont=*-italic,BoldItalicFont=*-bolditalic]
\setsansfont{texgyreheros}[Extension=.otf,UprightFont=*-regular,BoldFont=*-bold,ItalicFont=*-italic,BoldItalicFont=*-bolditalic]
\setmonofont{texgyrecursor}[Extension=.otf,UprightFont=*-regular,BoldFont=*-bold,ItalicFont=*-italic,BoldItalicFont=*-bolditalic,Scale=MatchLowercase]
\usepackage{amsmath,amssymb}
\usepackage{unicode-math}
\setmathfont{texgyrepagella-math.otf}
\usepackage{xcolor,microtype,enumitem,needspace,fancyhdr}
\usepackage{bookmark}
\definecolor{ink}{HTML}{17344B}
\definecolor{accent}{HTML}{197D80}
\definecolor{muted}{HTML}{596773}
\hypersetup{colorlinks=true,linkcolor=accent,urlcolor=accent,pdftitle={RA-11 - Optimal
trace-estimation complexity using Kronecker matrix-vector
queries},pdfauthor={Open Problems in Numerical Linear Algebra}}
\urlstyle{same}
\setlength{\parindent}{0pt}
\setlength{\parskip}{5pt plus 1pt minus 1pt}
\linespread{1.04}
\setlength{\emergencystretch}{3em}
\widowpenalty=10000
\clubpenalty=10000
\displaywidowpenalty=10000
\allowdisplaybreaks[1]
\setlist{nosep,leftmargin=1.5em}
\providecommand{\tightlist}{\setlength{\itemsep}{0pt}\setlength{\parskip}{0pt}}
\providecommand{\pandocbounded}[1]{#1}
\pagestyle{fancy}
\fancyhf{}
\fancyhead[L]{\sffamily\footnotesize\color{muted}OPEN PROBLEMS IN NUMERICAL LINEAR ALGEBRA}
\fancyhead[R]{\sffamily\bfseries\footnotesize\color{accent}RA-11}
\fancyfoot[L]{\parbox[b]{0.9\textwidth}{\sffamily\fontsize{8}{10}\selectfont\color{muted}Editorial ratings; a bounded search cannot certify that a problem remains unsolved.\\Literature check: 2026-09-14}}
\fancyfoot[R]{\sffamily\footnotesize\color{muted}\thepage}
\renewcommand{\headrulewidth}{0.3pt}
\renewcommand{\headrule}{\hbox to\headwidth{\color{accent}\leaders\hrule height \headrulewidth\hfill}}
\makeatletter
\renewcommand{\section}{\Needspace{5\baselineskip}\@startsection{section}{1}{0pt}{12pt plus 2pt minus 2pt}{4pt}{\sffamily\bfseries\large\color{ink}}}
\renewcommand{\subsection}{\Needspace{4\baselineskip}\@startsection{subsection}{2}{0pt}{9pt plus 1pt minus 1pt}{3pt}{\sffamily\bfseries\normalsize\color{ink}}}
\makeatother
\setcounter{secnumdepth}{0}
\begin{document}
{\sffamily\small\color{muted}Randomized and low-rank approximation\par}
\vspace{3pt}
{\raggedright\hyphenpenalty=10000\exhyphenpenalty=10000\sffamily\bfseries\fontsize{21}{24}\selectfont\color{ink}Optimal
trace-estimation complexity using Kronecker matrix-vector queries\par}
\vspace{7pt}
{\sffamily\small\textbf{Difficulty:} extreme\quad\textbf{Importance:} interesting
to the community\par}
{\sffamily\small\bfseries\color{accent}Status: Partially resolved\par}
\vspace{4pt}
{\color{accent}\hrule height 0.8pt}
\vspace{6pt}
\textbf{Rating rationale:} Extreme because optimal adaptive information
bounds remain unknown even across exponential tensor-order scales;
community impact is reliable computation with tensor-structured matrix
access.\\
\textbf{Topic:} Structured randomized trace estimation.

\subsection{Further partial results --- probability-sensitive probes, 14
September
2026}\label{further-partial-results-probability-sensitive-probes-14-september-2026}

\textbf{Author:} Sidney Holden, Center for Computational Biology,
Flatiron Institute, Simons Foundation.
\href{https://github.com/ajt60gaibb/OpenProblemsInNLA/blob/main/references/holden-further-2026-09-14/README.md}{Verified
affiliation and submission record}.

\href{https://github.com/ajt60gaibb/OpenProblemsInNLA/blob/main/references/holden-further-2026-09-14/RA-11/manuscript.pdf}{Theorems
2.1--2.2 and Corollary 2.3} give new upper bounds for every real PSD
input in the original full-vector real product-query model. In
particular, with \(N=n^q\),

\[Q(n,q,\varepsilon)\le\min\{N,2\lceil\sqrt{6N}/\varepsilon\rceil\}.\]

The constant \(6\) improves to \(3\) when \(N\) is a power of two. The
diagonal-control construction uses bounded-fair-bit, nonadaptive real
product queries. The fractional-moment branch has fixed-accuracy
exponential rate \(D_n=\log n+1-H_n\) in tensor order, where
\(H_n=\sum_{j=1}^n1/j\). Its sharpness is only for the specified
empirical-average estimator; it is not an unrestricted oracle lower
bound. A finite-bit version exists, with no efficient discretization
implementation claimed.

The new partial results passed a separate
\href{https://github.com/ajt60gaibb/OpenProblemsInNLA/blob/main/reviews/2026-09-14-further-submissions/RA-11-review.md}{independent
Codex AI-agent audit}. Exact diagonal and complex-response interpolation
checks and a floating moment grid passed; the numerical simulation suite
was not rerun. \textbf{The unrestricted joint minimax characterization
remains open}, so status remains Partially resolved. Inherited bounds in
the manuscript are outside this fresh audit; the prior independently
reviewed contribution below is retained.
\href{https://github.com/ajt60gaibb/OpenProblemsInNLA/blob/main/references/holden-further-2026-09-14/RA-11/manuscript.md}{Proof
text}. AI assistance is disclosed; no Lean verification, external human
peer review, formal verification or novelty claim is asserted.

\subsection{Partial results --- Sidney Holden, 13 September
2026}\label{partial-results-sidney-holden-13-september-2026}

\textbf{The general RA-11 target remains open.} Sidney Holden (Center
for Computational Biology, Flatiron Institute, Simons Foundation)
supplies
\href{https://github.com/ajt60gaibb/OpenProblemsInNLA/blob/main/references/holden-ra11-2026-09-13/manuscript/ra11_partial_results.pdf}{partial
results}, with
\href{https://github.com/ajt60gaibb/OpenProblemsInNLA/blob/main/references/holden-ra11-2026-09-13/manuscript/ra11_partial_results.tex}{proof
source},
\href{https://github.com/ajt60gaibb/OpenProblemsInNLA/blob/main/references/holden-ra11-2026-09-13/README.md}{verified
affiliation and submission record}, and a separate
\href{https://github.com/ajt60gaibb/OpenProblemsInNLA/blob/main/references/holden-ra11-2026-09-13/independent-review.md}{independent
informal Codex AI-agent audit}.

Theorem 1.1 gives general upper and lower bounds that do not match
throughout the parameter range. It yields
\(Q(n,q,\varepsilon)=\Theta(n^q)\) when \(\varepsilon\le n^{-q}\). Under
the additional promise that the unknown matrix is a tensor product of
PSD factors, the manuscript establishes
\(\Theta(\min\{n,\sqrt q/\varepsilon\})\) queries. It also gives a
\(q+1\)-call real simulation of a complex product query, promised
low-rank recovery, and a counterexample to the source's Conjecture 23.
That counterexample does not settle RA-11.

The audit passed these partial scopes; exact and numerical checks were
rerun. The surviving question is the unrestricted adaptive complexity
for arbitrary PSD matrices, uniformly in all three parameters. No full
resolution, external human peer review, novelty certification or formal
verification is claimed. No Lean verification was performed. The
original target follows unchanged.

Let \(n,q\ge2\) be integers and \(0<\varepsilon<1/2\). An unknown real
symmetric PSD matrix \(M\in\mathbb R^{n^q\times n^q}\) is accessible
only through an exact oracle which, on input
\(v_1,\ldots,v_q\in\mathbb R^n\), returns

\[M(v_1\otimes\cdots\otimes v_q)\in\mathbb R^{n^q}.\]

Let \(Q(n,q,\varepsilon)\) be the smallest integer \(t\) for which a
randomized algorithm can, for every such \(M\), make at most \(t\)
oracle calls and output \(\widehat t\in\mathbb R\) satisfying

\[\Pr\bigl\{|\widehat t-\mathop{\mathrm{tr}}\nolimits(M)|\le\varepsilon\mathop{\mathrm{tr}}\nolimits(M)\bigr\}\ge\frac23.\]

Queries may depend on previous responses and internal randomness. The
algorithm may perform arbitrary finite exact arithmetic between calls;
only oracle calls are counted. The inputs \(n,q,\varepsilon\) are known.
Neither the query vectors nor their concatenation are required to have
bounded condition number, and there is no restriction to Hutchinson
estimators.

Determine \(Q(n,q,\varepsilon)\) up to universal constant factors,
simultaneously in all three parameters.

This is a precise query-complexity formulation of the source's request
for tight trace-estimation bounds. It would quantify the cost of
exploiting rank-one tensor queries when ordinary matrix access is
unavailable.

\subsection{References}\label{references}

\begin{enumerate}
\def\labelenumi{\arabic{enumi}.}
\tightlist
\item
  R. A. Meyer, W. Swartworth, and D. P. Woodruff,
  \href{https://arxiv.org/html/2502.08029v2}{Understanding the Kronecker
  Matrix-Vector Complexity of Linear Algebra}, ICML 2025, PMLR 267,
  43909--43933. Definition 1 specifies the oracle; §6, final sentence,
  asks for tight trace-estimation bounds. Theorem 7 bounds a conditioned
  scalar quadratic-form oracle; its following discussion distinguishes
  full-vector responses.
  \href{https://proceedings.mlr.press/v267/meyer25a.html}{Published
  record}.
\item
  R. A. Meyer and H. Avron,
  \href{https://arxiv.org/abs/2309.04952}{Hutchinson's Estimator is Bad
  at Kronecker-Trace-Estimation}, \emph{SIAM Journal on Matrix Analysis
  and Applications} 47 (2026), 353--387, abstract and trace-estimator
  bounds. These tight rates concern the specified estimator and query
  distributions.
\end{enumerate}

\subsection{Status check --- 2026-09-08}\label{status-check-2026-09-08}

Checked the first paper's latest listed arXiv v2 (2025-02-13), §6, and
its ICML publication record; checked the second paper's latest listed v2
(2025-01-31) and 2026 journal record. Searches for ``Kronecker trace
estimation tight 2026'' and the exact first title found no solution for
unrestricted adaptive algorithms. The available trace lower bound
assumes both scalar quadratic-form queries and bounded query
conditioning; it does not establish a lower bound for the unrestricted
full-vector oracle defined here. The 2026 Hutchinson article establishes
estimator-specific rates and therefore does not determine \(Q\). The
minimax definition and explicit parameter range are editorial
formalization of a source-stated complexity question, not a quoted
formula from the paper.

\subsection{Audit --- 2026-09-10}\label{audit-2026-09-10}

Rechecked
\href{https://arxiv.org/html/2502.08029v2}{Meyer--Swartworth--Woodruff,
§6} and the \href{https://doi.org/10.1137/24M1720895}{2026 Hutchinson
publication}. Kronecker-trace-complexity searches found no unrestricted
adaptive characterization. Conditioned scalar-query lower bounds and
estimator-specific rates still do not settle the full-vector oracle.
\end{document}
